\documentclass[11pt]{article}
\usepackage[utf8]{inputenc}
\usepackage{amsmath,amssymb,amsthm}
\usepackage{url}
\usepackage{hyperref}

\newtheorem{theorem}{Theorem}
\newtheorem{lemma}[theorem]{Lemma}
\newtheorem{corollary}[theorem]{Corollary}
\newtheorem{conjecture}[theorem]{Conjecture}
\newtheorem{remark}[theorem]{Remark}
\theoremstyle{definition}
\newtheorem{definition}[theorem]{Definition}

\title{A Width Sufficient Condition for Graham's Distinct Partial Sums Conjecture}
\author{Autonomous Mathematical Research Project}
\date{\today}

\begin{document}

\maketitle

\begin{abstract}
Graham's conjecture (1971, repeated by Erd\H{o}s) asserts that every subset
$A\subseteq\mathbb F_p\setminus\{0\}$ admits an ordering whose partial sums
are pairwise distinct modulo $p$. The conjecture remains open in general; the
best-known results cover $|A|\le 12$, $|A|\in\{p-2,p-3\}$, and (recently)
$|A|\le \lceil \log p/\log\log p\rceil$ via rectification.

We prove a new, complementary sufficient condition. Lift $A$ to centered
representatives in $(-(p-1)/2,(p-1)/2)$, and let $P$ and $N$ be the positive
representatives and the magnitudes of the negative ones. If
$\sum P < p$ and $\sum N < p$, then $A$ is sequenceable
modulo $p$. This covers sets of size up to about $\sqrt{2p}$ (for instance
$A=\{1,\dots,k\}$ with $k(k+1)/2<p$), well beyond the rectification range.

The proof is elementary: a self-contained integer-ordering lemma, proved by
induction through repeatedly removing an element from the side of strictly
larger total, produces orderings whose prefix sums avoid each other; the
resulting integer partial sums are distinct and lie in an interval of width
$<p$, hence are distinct modulo $p$.

All claims are accompanied by explicit evidence categories; the theorem and
its proof are given unconditionally, and the construction is verified
computationally on all small cases and thousands of random instances. The
general conjecture remains open.
\end{abstract}

\section{Introduction}

\subsection{The problem}
For a finite subset $A$ of an abelian group, call an ordering
$a_1,\dots,a_n$ of its elements \emph{valid} if the partial sums
\[
  s_k=a_1+a_2+\cdots+a_k,\qquad 1\le k\le n,
\]
are pairwise distinct. Graham's conjecture (1971; repeated by Erd\H{o}s and
Graham \cite{EG80}, Problem 475) asserts:

\begin{conjecture}[Graham]
For every prime $p$ and every $A\subseteq\mathbb F_p\setminus\{0\}$, the set
$A$ admits a valid ordering modulo $p$.
\end{conjecture}

Such sets are called \emph{sequenceable}. The conjecture is connected to
sequencing of groups, complete mappings, Latin transversals (Snevily's
conjecture, now solved by Arsovski \cite{arsovski2011}), and the Hall--Paige
conjecture. Despite considerable attention \cite{costa2020, costa2022,
hicks2019}, it remains open in general.

\subsection{State of the art}
According to Kravitz \cite{kravitz2024}, the conjecture is known to hold for
$|A|\le 12$ (polynomial method / Combinatorial Nullstellensatz), for
$|A|\in\{p-2,p-3\}$ \cite{hicks2019}, and Kravitz proves it for
$|A|\le\lceil\log p/\log\log p\rceil$ using Lev's rectification theorem
\cite{lev2008} together with the integer version
\cite{kravitz2024,sawin2015}. The general case (arbitrary $|A|$) is open.

\subsection{Our result}
We give a sufficient condition that is orthogonal to cardinality: it depends
on the size of the centered representatives rather than on their number.

\begin{theorem}[Width sufficient condition]\label{thm:main}
Let $p$ be prime and $A\subseteq\mathbb F_p\setminus\{0\}$. Choose centered
representatives $a'\in\{-(p-1)/2,\dots,(p-1)/2\}\setminus\{0\}$ for each
$a\in A$, and put
\[
  P=\{a':a'>0\},\qquad N=\{-a':a'<0\}.
\]
If
\[
  \sum P < p,\qquad \sum N < p,
\]
then $A$ is sequenceable modulo $p$.
\end{theorem}

The condition includes sets of size up to roughly $\sqrt{2p}$. For example:

\begin{corollary}\label{cor:consecutive}
If $A=\{1,2,\dots,k\}\subseteq\mathbb F_p^\times$ satisfies
$k(k+1)/2<p$, then $A$ is sequenceable modulo $p$.
\end{corollary}

This is not implied by the rectification range when $k\gg\log p/\log\log p$.

\section{Definitions}

Let $A=\{a_1,\dots,a_n\}\subseteq\mathbb F_p\setminus\{0\}$. An ordering is
a permutation $b_1,\dots,b_n$ of $A$. Its partial sums are
$s_k=b_1+\cdots+b_k$. The ordering is valid if $s_1,\dots,s_n$ are pairwise
distinct in $\mathbb F_p$. Equivalently, no nonempty consecutive block sums
to zero. $A$ is sequenceable if some ordering is valid.

Centered representatives: each residue class mod $p$ is represented by its
unique element in $\{-(p-1)/2,\dots,(p-1)/2\}$.

\section{An integer ordering lemma}

\begin{lemma}\label{lem:ordering}
Let $P,N$ be finite sets of positive integers with $\sum P\ne\sum N$. Then
there exist orderings of $P$ and of $N$ such that no nonempty prefix-sum of
$P$ equals a prefix-sum of $N$ (the empty prefix-sum $0$ is common).
\end{lemma}

\begin{proof}
Induct on $|P|+|N|$. If one side is empty, its only prefix-sum is $0$ while
every nonempty prefix-sum of the other side is positive, so the claim holds.

Assume $\sum P>\sum N$ (the other case is symmetric). Then $P\ne\varnothing$.
Choose $p^*\in P$ with $p^*\ne\sum P-\sum N$: if $|P|\ge2$, at most one
element equals $\sum P-\sum N$, so another choice exists; if $P=\{p\}$ is a
singleton, then $p=\sum P>\sum N$ and $p\ne p-\sum N$ since $\sum N>0$.
Set $P'=P\setminus\{p^*\}$. Then $\sum P'=\sum P-p^*\ne\sum N$ by the
choice of $p^*$, so the induction hypothesis gives orderings of $P'$ and $N$
whose nonempty prefix-sums are disjoint.

Order $P$ as the ordering of $P'$ followed by $p^*$. Its nonempty prefix-sums
are those of $P'$ together with the full sum $\sum P$. Since
$\sum P>\sum N\ge$ (every nonempty prefix-sum of $N$), the full sum $\sum P$
is not a prefix-sum of $N$. Hence no nonempty prefix-sum of $P$ equals a
prefix-sum of $N$. The case $\sum N>\sum P$ is symmetric.
\end{proof}

\section{Proof of the main theorem}

\begin{proof}[Proof of Theorem \ref{thm:main}]
\textbf{Case 1:} $\sum P \ne \sum N$.

By Lemma 1 choose orderings of $P$ and $N$ whose nonempty prefix-sums are
disjoint. Output the ordering of $A'$
\[
  \operatorname{reverse}(P\text{-ordering})\quad\text{followed by}\quad
  (-n\ \text{for }n\in N\text{-ordering}).
\]
Let $S_1,\dots,S_n$ be the integer partial sums.

\textbf{Distinctness as integers.} The positive phase (reverse of the
$P$-ordering) has strictly increasing partial sums, all in $(0,\sum P]$.
The negative phase has strictly decreasing partial sums. A cross-phase
collision $S_i=S_j$ would require
\[
  \sum P-q=\sum P-n
\]
for $q$ a (proper or empty) prefix-sum of the $P$-ordering and $n$ a
nonempty prefix-sum of the $N$-ordering, i.e. $q=n$. By Lemma 1 this cannot
happen: no nonempty prefix-sum of $P$ equals a prefix-sum of $N$, and $q=0$
would require $n=0$, not allowed. Hence $S_1,\dots,S_n$ are distinct
integers.

\textbf{Range bound.} Every partial sum lies in $[m,M]$ with
$M=\sum P$ and $m=\min(1,\sum P-\sum N)$: the positive phase lies in
$(0,\sum P]\subseteq[1,\sum P]$, and the negative phase lies in
$[\sum P-\sum N,\sum P)$. Thus all partial sums lie in an interval of width
\[
  W=\sum P-\min(1,\sum P-\sum N).
\]
If $\sum P>\sum N$, then $\sum P-\sum N\ge1$, so $W=\sum P-1<p-1<p$.
If $\sum N>\sum P$, then $W=\sum N<p$. In either case $W<p$.

\textbf{Case 2:} $\sum P = \sum N$.

Since $A$ is a subset of nonzeros, if $P$ is empty then $N$ must be empty,
so $A$ is empty and sequenceable. Assume $P$ is nonempty. Choose any
$p^* \in P$ and let $P' = P \setminus \{p^*\}$. Then
$\sum P' = \sum P - p^* < \sum P = \sum N$. Since $P'$ and $N$ have unequal
sums, we can apply Lemma 1 to get orderings $P'\text{-ordering}$ and
$N\text{-ordering}$ whose nonempty prefix sums are disjoint.

Let the $P\text{-ordering}$ be the $P'\text{-ordering}$ followed by $p^*$.
The prefix sums of the $P\text{-ordering}$ are exactly those of the
$P'\text{-ordering}$, plus the full sum $\sum P$. Output the ordering
of $A'$
\[
  \operatorname{reverse}(P\text{-ordering})\quad\text{followed by}\quad
  (-n\ \text{for }n\in N\text{-ordering}).
\]
Let $S_1,\dots,S_n$ be the integer partial sums.

\textbf{Distinctness as integers.} Just as in Case 1, the positive and
negative phases are internally strictly monotone. A cross-phase collision
occurs iff $\sum P - q = \sum P - n$, iff $q=n$, where $q$ is a proper
prefix-sum of the $P\text{-ordering}$, and $n$ is a nonempty prefix-sum of
the $N\text{-ordering}$. Since $q$ must be a proper prefix sum of the
$P\text{-ordering}$, it is one of the prefix sums of the
$P'\text{-ordering}$. But by Lemma 1, no prefix sum of the
$P'\text{-ordering}$ equals a prefix sum of the $N\text{-ordering}$. Thus
$q=n$ is impossible, and there are no cross-phase collisions.

\textbf{Range bound.} The positive phase partial sums lie in $(0, \sum P]$.
The negative phase starts at $\sum P - \text{first element of } N\text{-ordering}$
and ends at $\sum P - \sum N = 0$. So all partial sums lie in $[0, \sum P]$.
The width is $W = \sum P < p$.

\textbf{Modulo $p$.} In both cases, the $S_k$ are distinct integers in an
interval of width $W<p$. If $S_i\equiv S_j\pmod p$ for $i\ne j$, then
$p\mid(S_i-S_j)$, so $|S_i-S_j|\ge p$, contradicting $|S_i-S_j|\le W<p$.
Hence the partial sums are pairwise distinct modulo $p$.
\end{proof}

\section{Limitations and remaining questions}

Graham's conjecture for arbitrary $|A|$ (in particular sets whose centered
representatives have large absolute sum) remains open. A natural strengthening
is the question of two-sided valid orderings raised in \cite{kravitz2024}.

\section{Evidence categories}

\begin{itemize}
\item \textbf{PROVED HERE:} Theorem \ref{thm:main}, Corollary
\ref{cor:consecutive}, Lemma \ref{lem:ordering}.
\item \textbf{ESTABLISHED IN PRIOR WORK:} the integer version of Graham's
conjecture \cite{kravitz2024,sawin2015}; Lev's rectification theorem
\cite{lev2008}; the state-of-the-art summary \cite{kravitz2024}.
\item \textbf{COMPUTATIONALLY VERIFIED (finite):} Theorem \ref{thm:main}'s
construction was checked exhaustively for every width-satisfying subset of
$\mathbb F_p^\times$ for $p\le 23$ and on thousands of random instances, with
complete agreement against brute force (see \texttt{tests/}, \texttt{src/}).
\end{itemize}

\section*{Acknowledgements}
This project was carried out autonomously. The problem was identified via the
Open Problem Garden and the arXiv preprint \cite{kravitz2024}; all code and
write-up are original to this project.

\begin{thebibliography}{99}

\bibitem{EG80} P. Erd\H{o}s and R. L. Graham, \emph{Old and new problems and
results in combinatorial number theory}, L'Enseignement Mathématique,
Université de Genève, 1980.

\bibitem{kravitz2024} N. Kravitz, Rearranging small sets for distinct partial
sums, arXiv:2407.01835 [math.CO], 2024.

\bibitem{sawin2015} W. Sawin, comment on ``Ordering subsets of the cyclic
group to give distinct partial sums,'' MathOverflow,
\url{https://mathoverflow.net/q/202857}, 2015.

\bibitem{lev2008} V. Lev, The rectifiability threshold in abelian groups,
\emph{Combinatorica} 28 (2008), 491--497.

\bibitem{blr1998} Y. Bilu, V. Lev, I. Ruzsa, Rectification principles in
additive number theory, \emph{Discrete Comput. Geom.} 19 (1998), 343--353.

\bibitem{hicks2019} J. Hicks, M. Ollis, J. Schmitt, Distinct partial sums in
cyclic groups: polynomial method and constructive approaches,
\emph{J. Combin. Des.} 27 (2019), 369--385.

\bibitem{costa2020} S. Costa, M. A. Pellegrini, Some new results about a
conjecture by Brian Alspach, \emph{Arch. Math. (Basel)} 115 (2020), 479--488.

\bibitem{costa2022} S. Costa, S. Della Fiore, M. Ollis, S. Rovner-Frydman,
On sequences in cyclic groups with distinct partial sums,
\emph{Electron. J. Combin.} 29 (2022), \#33.

\bibitem{arsovski2011} B. Arsovski, A proof of Snevily's conjecture,
\emph{Israel J. Math.} 182 (2011), 505--508.

\bibitem{hough2015} B. Hough, Solution of the minimum modulus problem for
covering systems, \emph{Ann. of Math.} 181 (2015), 361--382.

\end{thebibliography}

\end{document}
